\vspace{\presectionspace}
\section{Conclusions, Limitations and Societal Impact}
\label{sec:limitations}
\vspace{\postsectionspace}

\name showcases the effectiveness of frozen large pretrained language models as text encoders for the text-to-image generation using diffusion models. Our observation that scaling the size of these language models have significantly more impact than scaling the U-Net size on overall performance encourages future research directions on exploring even bigger language models as text encoders. Furthermore, through \name we re-emphasize the importance of classifier-free guidance, and we introduce dynamic thresholding, which allows usage of much higher guidance weights than seen in previous works. With these novel components, \name produces $1024\times1024$ samples with unprecedented photorealism and alignment with text.

Our primary aim with \name is to advance research on generative methods, using text-to-image synthesis as a test bed. While end-user applications of generative methods remain largely out of scope, we recognize the potential downstream applications of this research are varied and may impact society in complex ways. On the one hand, generative models have a great potential to complement, extend, and augment human creativity \cite{hughes2021}. Text-to-image generation models, in particular, have the potential to extend image-editing capabilities  and lead to the development of new tools for creative practitioners. On the other hand, generative methods can be leveraged for malicious purposes, including harassment and misinformation spread \cite{franks2019}, and raise many concerns regarding social and cultural exclusion and bias \cite{Srinivasan2021, Sequeira2021, Steed2021}. 
These considerations inform our decision to not to release code or a public demo. In future work we will explore a framework for responsible externalization that balances the value of external auditing with the risks of unrestricted open-access.

Another ethical challenge relates to the large scale data requirements of text-to-image models, which have have led researchers to rely heavily on large, mostly uncurated, web-scraped datasets. 
While this approach has enabled rapid algorithmic advances in recent years, datasets of this nature have been critiqued and contested along various ethical dimensions. For example, public and academic discourse regarding appropriate use of public data has raised concerns regarding data subject awareness and consent \cite{MegaPixels, dulhantychrisIssuesComputerVision2020, Scheuerman2021DoDH, Paullada2021}. Dataset audits have revealed these datasets tend to reflect social stereotypes, oppressive viewpoints, and derogatory, or otherwise harmful, associations to marginalized identity groups \cite{prabhu2020, birhane2021laionaudit}. 
Training text-to-image models on this data risks reproducing these associations and causing significant representational harm that would disproportionately impact individuals and communities already experiencing marginalization, discrimination and exclusion within society. As such, there are a multitude of data challenges that must be addressed before text-to-image models like \name can be safely integrated into user-facing applications. While we do not directly address these challenges in this work, an awareness of the limitations of our training data guide our decision not to release \name for public use. We strongly caution against the use text-to-image generation methods for any user-facing tools without close care and attention to the contents of the training dataset. 

\name's training data was drawn from several pre-existing datasets of image and English alt-text pairs. 
\ifdefined\neuripssubmission
A subset of this data was filtered to removed noise and undesirable content, such as pornographic imagery and toxic language. 
\else
400 million examples came from FIT400M\footnote{https://github.com/google-research/babeldraw/data_cards/fit400m.md}, a cleaned version of the internal Alt-Text dataset. This data was filtered to removed noise and undesirable content, such as pornographic imagery and toxic language. 
\fi
However, a recent audit of one of our data sources, LAION-400M \cite{schuhmann2021laion}, uncovered a wide range of inappropriate content including pornographic imagery, racist slurs, and harmful social stereotypes \cite{birhane2021laionaudit}. This finding informs our assessment that \name is not suitable for public use at this time and also demonstrates the value of rigorous dataset audits and comprehensive dataset documentation (e.g. \cite{gebruDatasheetsDatasets2020, datacards}) in informing consequent decisions about the model's appropriate and safe use. \name also relies on text encoders trained on uncurated web-scale data, and thus inherits the social biases and limitations of large language models \cite{bordia-naacl-2017, bender2021, weidinger2021}. 


 While we leave an in-depth empirical analysis of social and cultural biases encoded by \name to future work, our small scale internal assessments reveal several limitations that guide our decision not to release \name at this time.
 First, all generative models, including \name, 
 \name, may run into danger of dropping modes of the data distribution, which may further compound the social consequence of dataset bias. Second, \name exhibits serious limitations when generating images depicting people. Our human evaluations found \name obtains significantly higher preference rates when evaluated on images that do not portray people, indicating  a degradation in image fidelity. Finally, our preliminary assessment also suggests \name  encodes several social biases and stereotypes, including an overall bias towards generating images of people with lighter skin tones and a tendency for images portraying different professions to align with Western gender stereotypes. Even when we focus generations away from people, our preliminary analysis indicates \name encodes a range of social and cultural biases when generating images of activities, events, and objects. 

While there has been extensive work auditing image-to-text and image labeling models for forms of social bias (e.g. \cite{gendershades, Burns2018, Steed2021}), there has been comparatively less work on social bias evaluation methods for text-to-image models, with the recent exception of \citep{dalleval}. We believe this is a critical avenue for future research and we intend to explore benchmark evaluations for social and cultural bias in future work---for example, exploring whether it is possible to generalize the normalized pointwise mutual information metric \cite{npmi} to the measurement of biases in image generation models. There is also a great need to develop a conceptual vocabulary around potential harms of text-to-image models that could guide the development of evaluation metrics and inform responsible model release.  We aim to address these challenges in future work.
